\loadgeometry{narrow}%
\chapter*{第二瓶中之信}

\begingroup\kaiti%
我的使命已经结束。请不要悲伤，也不要怀念我。我的存在现已无益于世界的进程；毋宁说，我现在留下信息，也是为害。可是这个世界是多么美丽；我还想再多看一眼。可是这个世界是多么残酷；我还想再挽留一分。不过，现在任何言语已经无用。只有寄蜉蝣于天地，归羔羊之山野。

世界需要新的观测者，世界需要新的引导者；世界需要新的神明，世界自己的神明。

你知道世界的一切；你就是世界，世界就是你。这是分析命题，也是综合命题。你不知道的事根本不存在，我没办法告诉你不存在的事。既然如此，我还能告诉你什么？

不要听从柴郡猫的妄言。请你想起，你将回归自身，你将审判实存；你将前往天空的另一边，在世界的尽头建立你的国度。

\halfline\hfill%
雪莉
\endgroup\restoregeometry

\chapter{轨道酒吧}

“列车会驶向哪里呢？”铃乘坐在银河列车之上，感受着转向架冲击轨道带来的细碎震动。这个问题或许有了答案：前方就是有形的轨道的尽头。驶过最后一块枕木，列车却没有停下来的意思；它继续前行，驶入夜空；铃这才明白轨道和车轮不过是装饰物。“向下坠落的是过山车，而向上坠落的是银河列车。”铃作出令人满意的总结。窗外那些原本遥不可及的星辰，此刻仿佛被无形之手拉近，闪烁着十字形的光芒，就这样被列车碾过。喜鹊似乎还是不肯放弃冬季难得的温暖；它们伴着列车飞行，羽毛上密集而耀眼的光芒遮蔽了摇曳的星光。“看来‘落霞与孤鹜齐飞’不如‘升斗与群鸦竞矗’。”铃思考着意象的选取，“那么下一句就该是‘冻星共江雪同光’了。”然而南方是不会下雪的；那么这样的诗句不过是铃的矫作。铃放弃了写诗的念头。她望着窗外飞快淌过的银河，忍不住哼起歌来：
\begin{quote}
跨越银河铁道之夜\\
终将与你再次相会\\
我们邀请星辰共舞\\
我们期翼春风同浴\par
\halfline%
我孑然一身\halfspace 立于天穹之上\\
请牵起我的手\\
请赐予我一束丹红\par
\halfline%
我们互相的话语\\
如今才回想起来\\
响彻此刻的夜中\par
\end{quote}

没有伴奏的歌唱是很滑稽的，尤其当你是唱的那个人。但是彩名还是表达了赞许：“很好听哦！铃可以当少女乐队的主唱了！”那么显然还需要吉他手和键盘手；贝斯可以先忽略。铃不知道坐在一旁的画匠会不会弹电子琴；希望他也能把摇滚视为一种他认可的艺术。他只是伸出带着老茧的手指在桌面上敲击着，仿佛在尝试为铃的唱词配上鼓点。

“真聒噪！”彩名对背景中的不和谐音感到不满；它们显然不是来为铃伴奏的。窗外的喜鹊渐渐跟不上持续加速的列车；它们的轨迹划出向后的流线，仅仅留下残存的叽喳。彩名打开车窗，对着鸟儿下达命令：“别跟来了！”她接着举起她的长伞试图击落鸟儿，却被它们轻盈地躲避。“你在干什么？”铃对彩名怪诞的行为表示不解。“它们已经完成了使命。”彩名说着铃听不懂的话语。

“我们到达目的地了。”画匠突然发言；然而列车并没有停下的意思。

“轨道酒吧自然需要在轨道上；它是一个移动的酒吧。如若不然，它便与普通的酒吧无异。”画匠对着两张怀疑的脸庞解释；他邀请她们前往后方的车厢。

两人跟随着画匠，通过车厢的连接处；自动门抱怨似的开始了它的工作，伴随着疲劳的声响缓缓开启。铃看向一旁的双针压力表；制动风管的气压是令人惊讶地正常。“原来不是逆汽制动，”铃对这种混搭感到奇怪，“所以使用蒸汽机车完全是出于文学目的。”她得出结论。

门后是另外一个世界：炫彩的灯带缠绕在各处，不与任何人类能够设想的扭结同痕；悬在空中的球面镜则反射着灯带的光线，将酒吧的空间分为不同颜色的离散区块。熙攘的身影在欢快的乐曲中扭动着身体；铃对乐曲大杂烩般的配器感到不满。不过好在她们无需参与其中；画匠带领她们走向一个相比之下稍微安静的角落，在那里是熟悉的身影。

“莉莉姆！艾莉丝！”臭蛤蟆举起一只蹼向她们招呼，“竟然能在这里见到你们。”彩名赶紧向他递眼色。

“我是铃；你怎么总是认错。”铃纠正他的错误；她和艾莉丝不是双胞胎啊。

“抱歉抱歉。我记混了。”臭蛤蟆挠了挠他的秃头，“诸位请坐。我们正举办茶会呢！”酒吧显然不是举办茶会的合适场所。铃和彩名以及画匠便在这张劣质木头制成的八仙桌旁落了座。

“哈，黑烟！”桌旁的另一位参与者，带着硕大礼帽的帽匠先前紧紧地盯着窗外，此时则向后靠到椅背上，高兴地叫出了声，“我赢了！现在你们都得喝。”他开始调整他的兴奋的帽子，“包括新来的各位！”

“事先说好：”画匠提醒似的开口，“我对你们的赌局毫无兴趣；你们别打扰到我的创作。”他取出速写本和炭笔；看来艺术家就需要把创作工具随身携带。

“‘创作’！”帽匠嘲讽道，“我看不如叫‘射精’。”

“你怎么能在两位小姐面前说出如此污秽之语，”臭蛤蟆赶紧伸出蹼挡住帽匠的嘴巴，“别玩你那一套了；现代心理学已经把它淘汰了。”接着，臭蛤蟆用头顶的大眼睛替帽匠向铃和彩名道歉。

“我又没在讨论心理学。你个只会按照诊断标准给人贴标签的蠢货……”被捂嘴的帽匠仍然在发表他的纯粹人身攻击。“总之，你们都得喝！”他拿起茶壶给在座的每个人斟满了一杯。铃望向面前的茶盏，里面是粘稠的乌黑。在酒吧的茶会上饮用咖啡，这即否定了酒吧的意义，又否定了茶会的确定性。

“老马已经喝醉了；你别给他倒了。”臭蛤蟆发表延迟的抱怨。老马趴在八仙桌的一角，长长的一侧脸颊摊在桌面上；他的巨大的鼻腔发出骇人的鼾声。“原来马除了站立着睡觉，也是可以趴下睡觉的。”铃记下这个事实。

彩名高兴地往咖啡中撒盐；铃不知道她手上的盐罐是从哪里来的；可能是酒吧提供的吧。一粒粒纯白色的晶体如暴风雪没入乌黑的湖面，泛起绵密的气泡，快要流溢出杯际。彩名大口地喝下这杯一定过咸的咖啡，似乎含在口中嚼了嚼，牙齿抱怨似的发出咯吱的声响。相比之下，铃抗拒喝这种看上去就不好喝的饮品。她尝试转移话题以暗度陈仓：“你们在赌什么？”

“我们刚才在赌上一轮投票能否选出神明；”帽匠为这场荒谬的赌博提供一个确定的意义，“接下来我们要赌下一轮投票能否选出神明。”听他的语气更像是在选举小组长，“这很简单；你看远处的那边，如果升起黑烟说明此轮投票无效，如果升起白烟说明世界有了新的神明。”

“神明是通过选举产生的吗？”铃质疑这样做的合法性，“如果成为神明需要选举作为根据，那么神明就不是完满的存在。”

“什么完满的存在。”帽匠似乎在怀疑铃是否是从中世纪穿越到现代的，“我们不过是在选举一位对象a。”

“还不如叫‘嫌疑人X’！”彩名插嘴；或许她推理小说看太多了。

“这是个好的类比。”帽匠竟然表达了认可，“我要好好想想追寻真凶是否是侦探的享乐。”他从帽子里掏出一个怪异的终端；这是他的AI助手。与其说是“他”在想，不如说是他在让大语言模型帮他想。“这是我们新研发的‘哲学模型’。”帽匠为它指定了一个名字。

“别摆弄你的赛博阳具了。”臭蛤蟆用帽匠的方式戏谑着；但他立刻感到后悔，大眼珠子向铃和彩名的方向瞟了瞟。铃示意她觉得无所谓；毕竟都是成年人了。至于彩名，魔法少女的表象世界里应该不存在这种东西。

“莉莉姆的手指会更舒服！”彩名脱线地展示着她刚在泡温泉前剪完指甲的手指。铃赶紧捂住彩名的嘴；千万别把她俩的关系透露出去。

帽匠重新被对话吸引注意。他做出总结：“确实。一看到手指，立刻想到能指，立刻想到菲勒斯，立刻想到生殖器，立刻想到享乐的错位，立刻想到符号的强奸。”中国人的想象唯在这一层能够如此滑动；铃为他补上这一句。“这样看来，LLM确实可以说是公用的阳具。”
